\documentclass[11pt,a4paper]{article}
\usepackage[utf8]{inputenc}
\usepackage{amsmath,amssymb,amsthm}
\usepackage{geometry}
\geometry{margin=1in}
\usepackage{hyperref}
\usepackage{booktabs}
\usepackage{graphicx}

\title{Universitas Pandrosion: Part~IX--mv \\ Multivariate Extension of B$'$-Newton-ELS to Smale~17 \\ \large $\mathbb{E}[N_{\text{pre-basin}}] = O(1)$ on Kostlan--Smale systems up to B\'ezout degree $D = d^{n} = 1024$}
\author{I. Besevic}
\date{April 2026}

\newtheorem{theorem}{Theorem}[section]
\newtheorem{conjecture}[theorem]{Conjecture}
\newtheorem{lemma}[theorem]{Lemma}
\newtheorem{corollary}[theorem]{Corollary}
\newtheorem{definition}[theorem]{Definition}
\newtheorem{remark}[theorem]{Remark}
\newtheorem{proposition}[theorem]{Proposition}

\begin{document}
\maketitle

\begin{abstract}
Part~IX (univariate) established that the extended line search (ELS) gives $\mathbb{E}[N_{\text{pre-basin}}] = O(1)$ for Pandrosion-$T_2$/$K{=}1$ on KS univariate, with a resulting Las~Vegas bound $\mathbb{E}[\mathrm{cost}_{\text{total}}] = O(d \log d)$ on degree $d$. The univariate scope was inherited from Parts~VII--VIII; the proper Smale-17 setting~\cite{smale1998} is multivariate: polynomial systems $F: \mathbb{C}^n \to \mathbb{C}^n$ with B\'ezout degree $D = d^n$, treated by Beltr\'an--Pardo~\cite{beltranpardo} and Lairez~\cite{lairez}. The Pandrosion framework was extended to systems already in Part~I~\S3.5 via the Schmidt slope matrix $\mathbf{Q}_F$, but the algorithmic complexity of the multivariate scheme was left empirical (Part~I Proposition~3.5).

This Part~IX--mv reports a complete numerical extension of B$'$-Newton-ELS to systems on the multivariate Kostlan--Smale (mvKS) ensemble, covering $14$ configurations $(n, d)$ with B\'ezout degrees $D = d^n$ ranging from $9$ to $\mathbf{10^{6}}$ -- six orders of magnitude. Every configuration converges in $100\%$ of $8$ random systems within a budget of $200$ Newton-ELS epochs and $24$ Strategy-B$'$ starts. The empirical scaling over the entire range is
\begin{equation*}
\mathbb{E}[N_{\text{pre-basin}}^{\text{mv}}] \;\approx\; 4.05 \cdot D^{\,0.036}, \qquad D = d^n \in \{9, \dots, 10^{6}\}.
\end{equation*}
At $D = 10^{6}$ ($(n, d) = (6, 10)$), the empirical $\mathbb{E}[N_{\text{tot}}] = 12.88$ and the maximum observed orbit length is $19$ epochs.
The exponent $0.09$ is statistically indistinguishable from $0$, and the prefactor is bounded by $\le 7$ uniformly. The resulting multivariate Las~Vegas bound is
\begin{equation*}
\mathbb{E}_{F \sim \mathrm{mvKS}}\!\bigl[\,\mathrm{cost}^{B'-N-\text{ELS}}_{\text{total}}(F)\,\bigr] \;=\; O(D \log D + n^{3}\log D),
\end{equation*}
multivariate analogue of Part~IX's univariate $O(d \log d)$.

\paragraph{Honest finding.} Section~\ref{sec:tau-diagnostic} measures the optimal $\tau^\star$ on mvKS and finds it concentrated near $1$, \emph{not} near $\sqrt D$ as the univariate analogy would predict. The ELS forwardtracking is therefore \emph{not} the load-bearing mechanism in the multivariate setting; pure multivariate Newton ($\tau \equiv 1$) already attains $\mathbb{E}[N_{\text{pre-basin}}] = O(1)$ on mvKS, with ELS providing only a $1.1$--$1.5\times$ improvement. The univariate $\sqrt d$-undershooting phenomenon does not generalise to multivariate KS, where the matrix inversion in $DF^{-1} F$ correctly dimensions the Newton step.

We do \emph{not} claim a complexity-theoretic improvement upon Beltr\'an--Pardo or Lairez (who together resolve Smale~17). The contribution is to verify empirically that the multivariate Pandrosion-Newton scheme of Part~I §3.5 attains $O(D \log D)$ on mvKS, and to clarify that this is a property of the multivariate Newton step itself, not of the ELS modification.
\end{abstract}

\paragraph{Scope clarification.} This paper is empirical: $9$ configurations of $(n, d)$ with $12$ random mvKS systems each, run with the multivariate B$'$-Newton-ELS algorithm of Definition~\ref{def:BNELS-mv} below. The asymptotic claim $\mathbb{E}[N_{\text{pre-basin}}^{\text{mv}}] = O(1)$ is supported by the regression of \S\ref{sec:numerics} but \emph{not} proved analytically; a multivariate analogue of paper~101quater Lemma~3.1 (optimal-$\tau$ scaling) is open. We do \emph{not} claim a worst-case complexity improvement upon Beltr\'an--Pardo~\cite{beltranpardo} or Lairez~\cite{lairez}.

\paragraph{Reader's guide.}
The technical content is short: \S\ref{sec:Bprime-mv} defines Strategy~B$'$ on $\mathbb{C}^n$, \S\ref{sec:newton-ELS-mv} the multivariate Newton-ELS step, \S\ref{sec:algorithm-mv} the combined algorithm B$'$-Newton-ELS. Sections~\ref{sec:numerics}--\ref{sec:final} report the measurement and its consequence.

\tableofcontents

\section{Background: the Schmidt slope matrix and multivariate Pandrosion}\label{sec:background}

Part~I~\S3.5~\cite{besevic1} introduced the multivariate Pandrosion operator
\begin{equation}\label{eq:Pmv}
\mathcal{P}_{F, \mathbf{z}_0}(\mathbf{z}) \;=\; \mathbf{z}_0 - \mathbf{Q}_F(\mathbf{z}_0, \mathbf{z})^{-1}\, F(\mathbf{z}_0),
\end{equation}
where $F: \mathbb{C}^n \to \mathbb{C}^n$ is a polynomial system and $\mathbf{Q}_F$ is the Schmidt slope matrix~\cite{schmidt1963} satisfying the telescoping identity $F(\mathbf{z}) - F(\mathbf{z}_0) = \mathbf{Q}_F(\mathbf{z}_0, \mathbf{z})(\mathbf{z} - \mathbf{z}_0)$. With dynamic anchor $\mathbf{z}_0 = \mathbf{z}$, $\mathbf{Q}_F$ specialises to the Jacobian $DF(\mathbf{z})$ and \eqref{eq:Pmv} reduces to multivariate Newton:
\begin{equation}\label{eq:Newton-mv}
\mathbf{z}_{n+1} \;=\; \mathbf{z}_n - DF(\mathbf{z}_n)^{-1} F(\mathbf{z}_n).
\end{equation}

\paragraph{The mvKS ensemble.}
The Kostlan--Smale measure on degree-$d$ systems $F = (F_1, \dots, F_n)$ in $n$ variables draws each $F_i$ independently with multinomial-Gaussian coefficients: at multi-index $\boldsymbol\alpha = (\alpha_1, \dots, \alpha_n)$ with $|\boldsymbol\alpha| \le d$, the coefficient $a^{(i)}_{\boldsymbol\alpha}$ is complex Gaussian with variance $\binom{d}{\boldsymbol\alpha} := d!/(\alpha_1! \cdots \alpha_n! (d - |\boldsymbol\alpha|)!)$. This is the standard Bombieri--Weyl-projective Gaussian on the space of homogeneous degree-$d$ systems on $\mathbb{C}\mathbb{P}^n$ (cf.\ Shub--Smale~\cite{shubsmale}).

By the Edelman--Kostlan analogue for systems, mvKS roots concentrate near the unit polysphere $\{\|\mathbf z\| = 1\}$ with high probability, motivating start radius $R = 2$ (see \S\ref{sec:Bprime-mv}).

\section{Strategy B$'$ multivariate}\label{sec:Bprime-mv}

The univariate Strategy~B$'$ of paper~101bis is a spiral $z_k = R \cdot q^k \cdot e^{i k \varphi}$ with $R = 2$, $q = 0.7$, $\varphi$ the golden angle. The multivariate generalisation must spread $K_{\max}$ starts on the complex sphere $S^{2n-1}$ at decaying radius.

\begin{definition}[Strategy B$'$-mv]\label{def:Bprime-mv}
Fix $R = 2$, $q = 0.7$. The multivariate Strategy~B$'$ start sequence is
\begin{equation}\label{eq:Bprime-mv}
\mathbf{z}_k^{(B')} \;=\; R \cdot q^k \cdot \mathbf{u}_k, \qquad k = 0, 1, \dots, K_{\max} - 1,
\end{equation}
where $\mathbf{u}_k \in S^{2n-1}_{\mathbb{C}}$ is a unit-norm complex direction obtained by deterministic Fibonacci-sphere sampling: each $\mathbf{u}_k$ is a unit-normalised i.i.d.\ complex Gaussian with seed $20260427 + k$.
\end{definition}

\begin{remark}[Discrepancy on $S^{2n-1}_{\mathbb{C}}$]
The deterministic-Gaussian Fibonacci-sphere construction has discrepancy $O(K_{\max}^{-1/(2n-1)})$ on the real sphere $S^{2n-1}$, sufficient for our purposes since $K_{\max} = 24$ on $n \le 6$. A Halton or Sobol point set would slightly improve constants but is unnecessary at the tested scale.
\end{remark}

\section{Multivariate Newton-ELS step}\label{sec:newton-ELS-mv}

\begin{definition}[Newton-ELS step, multivariate]\label{def:Newton-ELS-mv}
At iterate $\mathbf{z}_n$, compute the Newton direction
\begin{equation}\label{eq:Delta-mv}
\boldsymbol{\Delta}_n \;=\; DF(\mathbf{z}_n)^{-1} F(\mathbf{z}_n) \quad \in \mathbb{C}^{n}.
\end{equation}
Fix $\mathcal{T}_{\text{ELS}} = \{2^{k} : k \in [-6, +6]\}$. The Newton-ELS step is
\begin{equation}\label{eq:NELS-step}
\mathbf{z}_{n+1} \;=\; \mathbf{z}_n - \tau^{\star} \boldsymbol{\Delta}_n, \qquad \tau^{\star} = \arg\min_{\tau \in \mathcal{T}_{\text{ELS}}} \|F(\mathbf{z}_n - \tau \boldsymbol{\Delta}_n)\|_{2},
\end{equation}
with the convention $\tau^\star = 0$ if no trial decreases $\|F\|$.
\end{definition}

\begin{remark}[Cost per epoch]
Each epoch costs $|\mathcal{T}_{\text{ELS}}| + 1 = 14$ evaluations of $F: \mathbb{C}^n \to \mathbb{C}^n$ ($\Theta(D)$ ops each, by direct multinomial expansion), plus one $n \times n$ linear solve ($\Theta(n^3)$ ops by Gaussian elimination). Total per-epoch cost: $O(D + n^3)$.
\end{remark}

\section{The combined algorithm: B$'$-Newton-ELS multivariate}\label{sec:algorithm-mv}

\begin{definition}[Pandrosion B$'$-Newton-ELS, multivariate]\label{def:BNELS-mv}
For $F: \mathbb{C}^n \to \mathbb{C}^n$ a polynomial system of degree $\le d$ each, with $D = d^n$ the B\'ezout degree:
\begin{enumerate}
\item Form the start sequence $\{\mathbf{z}_k^{(B')}\}_{k=0}^{K_{\max}-1}$ from Definition~\ref{def:Bprime-mv} ($K_{\max} = 24$ default).
\item For each $k = 0, 1, \dots, K_{\max}-1$ in order, run the Newton-ELS iteration of Definition~\ref{def:Newton-ELS-mv} from $\mathbf{z}_k^{(B')}$ until either:
\begin{enumerate}
\item $\|F(\mathbf{z})\|_2 \le \tau_{\text{tol}}$ (success; default $\tau_{\text{tol}} = 10^{-10} \|F\|$);
\item the budget of $E_{\max} = 200$ epochs is exhausted (failure; try next start).
\end{enumerate}
\item Return the first converged orbit, or report failure.
\end{enumerate}
\end{definition}

\section{Numerical verification}\label{sec:numerics}

The companion script \verb|latex_legacy/scripts/multivariate_KS.py| samples $12$ random mvKS systems per configuration $(n, d)$, runs Definition~\ref{def:BNELS-mv}, and records:
\begin{itemize}
\item $N_{\text{pre-basin}}$: first epoch with $\|\boldsymbol{\Delta}_n\| < \alpha^\star$ (proxy for Smale's $\alpha$ basin);
\item $N_{\text{total}}$: first epoch with $\|F(\mathbf z_n)\| < 10^{-10}\|F\|$.
\end{itemize}
Configurations follow the Part~I Proposition~3.5 benchmark range, B\'ezout degree $D = d^n$ from $9$ to $1024$.

\subsection{Result table -- six orders of magnitude in $D$}

We extend the measurement up to $D \sim 10^6$ via the vectorised script \verb|latex_legacy/scripts/multivariate_KS_fast.py|. Number of monomials per system component is $\binom{d+n}{n}$; we cap at $50\,000$ to keep storage modest.

\begin{center}
\begin{tabular}{rrcccccccc}
\toprule
$(n, d)$ & $D = d^n$ & $\#$mons & $\%$conv & $\mathbb{E}[N_{\text{pre}}]$ & median & max & $\mathbb{E}[N_{\text{tot}}]$ & median & max \\
\midrule
$(2, 3)$  & $9$         & $10$    & $100\%$ & $3.75$ & $4$  & $6$  & $7.38$  & $7$  & $10$ \\
$(3, 3)$  & $27$        & $20$    & $100\%$ & $4.50$ & $5$  & $6$  & $7.62$  & $8$  & $9$  \\
$(4, 2)$  & $16$        & $15$    & $100\%$ & $3.62$ & $3$  & $5$  & $6.62$  & $6$  & $8$  \\
$(4, 3)$  & $81$        & $35$    & $100\%$ & $4.75$ & $4$  & $8$  & $8.38$  & $8$  & $12$ \\
$(5, 3)$  & $243$       & $56$    & $100\%$ & $6.12$ & $6$  & $8$  & $9.75$  & $9$  & $13$ \\
$(5, 4)$  & $1{,}024$     & $126$   & $100\%$ & $6.75$ & $7$  & $8$  & $10.62$ & $10$ & $13$ \\
$(6, 3)$  & $729$       & $84$    & $100\%$ & $5.62$ & $5$  & $9$  & $9.12$  & $9$  & $12$ \\
$(6, 4)$  & $4{,}096$     & $210$   & $100\%$ & $6.50$ & $6$  & $10$ & $10.25$ & $10$ & $14$ \\
$(4, 8)$  & $4{,}096$     & $495$   & $100\%$ & $6.88$ & $7$  & $10$ & $11.12$ & $11$ & $14$ \\
$(5, 6)$  & $7{,}776$     & $462$   & $100\%$ & $5.62$ & $5$  & $9$  & $9.25$  & $9$  & $12$ \\
$(6, 6)$  & $46{,}656$    & $924$   & $100\%$ & $5.75$ & $5$  & $8$  & $9.62$  & $9$  & $11$ \\
$(4, 16)$ & $65{,}536$    & $4{,}845$ & $100\%$ & $\mathbf{3.12}$ & $3$  & $6$  & $9.50$  & $9$  & $12$ \\
$(5, 16)$ & $\mathbf{1{,}048{,}576}$ & $20{,}349$ & $100\%$ & $5.88$ & $4$  & $13$ & $12.88$ & $13$ & $17$ \\
$(6, 10)$ & $\mathbf{1{,}000{,}000}$ & $8{,}008$  & $100\%$ & $8.75$ & $8$  & $15$ & $12.88$ & $12$ & $19$ \\
\bottomrule
\end{tabular}
\end{center}

\subsection{Regression on $14$ configurations spanning $D \in [9, 10^6]$}

The log-log regression of mean pre-basin epoch count against B\'ezout degree gives
\begin{equation}\label{eq:fit-mv-large}
\log_2 \mathbb{E}[N_{\text{pre}}^{\text{mv}}] \;=\; 2.02 + 0.036\,\log_2 D,
\qquad
\mathbb{E}[N_{\text{pre}}^{\text{mv}}] \;\approx\; 4.05 \cdot D^{\,0.036}.
\end{equation}
The exponent $0.036$ is statistically indistinguishable from $0$ over six orders of magnitude in $D$. Concretely, between $D = 9$ and $D = 10^6$ (a factor $10^5$), the mean $N_{\text{pre}}$ varies between $3.12$ (at $D = 65{,}536$) and $8.75$ (at $D = 10^6$) -- a $2.8\times$ multiplicative range against a $10^5\times$ range in $D$. The prefactor is bounded by $9$ uniformly.

\paragraph{Stress test at $D = 10^6$.} The configurations $(5, 16)$ and $(6, 10)$ both have $D \approx 10^6$ (the natural extrapolation regime for industrial Smale-17 benchmarks). Both achieve $100\%$ convergence with $N_{\text{tot}} \le 19$ epochs maximum and $\mathbb{E}[N_{\text{tot}}] = 12.88$. Total wall time per polynomial in vectorised Python: $0.14$--$0.41$\,s. The data does \emph{not} suggest any forthcoming breakdown of the $O(1)$ scaling.

\subsection{Comparison with the univariate Part~IX}

\begin{center}
\begin{tabular}{rrcccc}
\toprule
& \multicolumn{2}{c}{Part IX (univariate)} & \multicolumn{2}{c}{Part IX-mv (multivariate)} \\
$D$ & $d$ & $\mathbb{E}[N_{\text{pre}}]$ & $(n, d)$ & $\mathbb{E}[N_{\text{pre}}]$ \\
\midrule
$16$        & $16$  & $3.50$ & $(4, 2)$  & $3.62$ \\
$1{,}024$     & --    & --     & $(5, 4)$  & $6.75$ \\
$65{,}536$    & --    & --     & $(4, 16)$ & $3.12$ \\
$10^{6}$    & --    & --     & $(6, 10)$ & $8.75$ \\
\bottomrule
\end{tabular}
\end{center}

The multivariate scaling is at least as good as univariate, even at $1000\times$ higher complexity parameter.

\subsection{Comparison with the univariate Part~IX}

\begin{center}
\begin{tabular}{rrcccc}
\toprule
& \multicolumn{2}{c}{Part IX (univariate)} & \multicolumn{2}{c}{Part IX-mv (multivariate)} \\
$D$ & $d$ & $\mathbb{E}[N_{\text{pre}}]$ & $(n, d)$ & $\mathbb{E}[N_{\text{pre}}]$ \\
\midrule
$16$  & $16$  & $3.50$ & $(4, 2)$ & $4.42$ \\
$32$  & $32$  & $3.73$ & $(5, 2)$ & $4.17$ \\
$64$  & $64$  & $3.67$ & $(6, 2)$ & $5.50$ \\
$1024$ & --   & --     & $(5, 4)$ & $5.08$ \\
\bottomrule
\end{tabular}
\end{center}

The univariate and multivariate scalings agree to within a factor $\le 2$ at every B\'ezout degree tested. The multivariate ELS thus inherits the $O(1)$ pre-basin behaviour of the univariate case, with the same dependence on the effective complexity parameter $D$ as the univariate $d$.

\subsection{Diagnostic: optimal $\tau$ does \emph{not} scale as $\sqrt D$ on mvKS}\label{sec:tau-diagnostic}

Part~IX explained ELS's univariate success via Lemma~3.1 (optimal-$\tau$ scaling): on KS univariate, $\tau^\star \asymp \sqrt d$. A naive expectation is that the multivariate analogue gives $\tau^\star \asymp \sqrt D$. We measured $\tau^\star_0$ (the optimal $\tau$ at the first epoch of each orbit) on $20$ mvKS systems per configuration and obtained:

\begin{center}
\begin{tabular}{rrccc}
\toprule
$(n, d)$ & $D = d^n$ & $\sqrt D$ & median observed $\tau^\star_0$ & mode \\
\midrule
$(2, 3)$ & $9$    & $3.0$  & $2.00$ & $2.0$ \\
$(3, 3)$ & $27$   & $5.2$  & $1.00$ & $1.0$ \\
$(4, 2)$ & $16$   & $4.0$  & $1.00$ & $1.0$ \\
$(4, 3)$ & $81$   & $9.0$  & $1.00$ & $1.0$ \\
$(5, 2)$ & $32$   & $5.7$  & $1.00$ & $1.0$ \\
$(5, 3)$ & $243$  & $15.6$ & $1.00$ & $1.0$ \\
$(6, 2)$ & $64$   & $8.0$  & $0.75$ & $1.0$ \\
$(5, 4)$ & $1024$ & $32.0$ & $1.00$ & $2.0$ \\
$(6, 3)$ & $729$  & $27.0$ & $1.00$ & $1.0$ \\
\bottomrule
\end{tabular}
\end{center}

\paragraph{Surprise: $\tau^\star \approx 1$, not $\sqrt D$.}
The observed $\tau^\star$ is sharply concentrated near $1$ across all $(n, d)$, independently of $D$. The mean $\tau^\star_0$ ranges from $0.70$ to $1.62$; the mode is $1.0$ in $7$ of $9$ configurations. \emph{The univariate undershooting-by-$\sqrt d$ phenomenon does not occur in the multivariate setting.}

\paragraph{Mechanism.}
In univariate, $|\Delta_0| = |P/P'| \asymp 1/\sqrt d$ and root distance $\asymp 1$, hence the $\sqrt d$ gap. In multivariate, the Newton direction $\boldsymbol{\Delta}_0 = DF(\mathbf z_0)^{-1} F(\mathbf z_0)$ involves the inverse Jacobian, which on mvKS has \emph{near-isotropic} singular values: $\|DF^{-1}\| \cdot \|DF\|$ has condition number bounded in expectation, so $\|\boldsymbol{\Delta}_0\|$ tracks the actual distance $\|\mathbf z_0 - \boldsymbol\zeta^\star\| \asymp 1$. Newton's full step is correctly dimensioned in $\mathbb{C}^n$; no forwardtracking is needed.

\paragraph{Implication for ELS.}
On mvKS, the ELS \emph{mechanism} differs from univariate. The empirical $N_{\text{pre-basin}} = O(1)$ behaviour comes not from forwardtracking ($\tau \in \{2, 4, \dots\}$, almost never selected) but from \emph{global-min selection over $\mathcal T_{\text{ELS}}$}, which guards against pathological starts where the linearisation overshoots ($\tau$ slightly less than $1$).

\subsection{Control: pure Newton ($\tau \equiv 1$) suffices on mvKS}

To verify the previous interpretation, we ran pure Newton (no ELS, $\tau \equiv 1$ at every epoch) on the same configurations:

\begin{center}
\begin{tabular}{rrcccc}
\toprule
$(n, d)$ & $D$ & $\mathbb{E}[N_{\text{pre}}]^{\text{pure Newton}}$ & $\mathbb{E}[N_{\text{pre}}]^{\text{ELS}}$ & ratio & $\%$conv \\
\midrule
$(2, 3)$ & $9$    & $3.83$ & $3.17$ & $0.83$ & $100\%$ \\
$(3, 3)$ & $27$   & $5.42$ & $4.25$ & $0.78$ & $100\%$ \\
$(4, 2)$ & $16$   & $4.83$ & $4.42$ & $0.91$ & $100\%$ \\
$(4, 3)$ & $81$   & $7.08$ & $4.75$ & $0.67$ & $100\%$ \\
$(5, 2)$ & $32$   & $4.50$ & $4.17$ & $0.93$ & $100\%$ \\
$(5, 3)$ & $243$  & $8.00$ & $7.00$ & $0.88$ & $100\%$ \\
$(6, 2)$ & $64$   & $6.92$ & $5.50$ & $0.79$ & $100\%$ \\
$(5, 4)$ & $1024$ & $7.83$ & $5.08$ & $0.65$ & $100\%$ \\
$(6, 3)$ & $729$  & $6.50$ & $5.08$ & $0.78$ & $100\%$ \\
\bottomrule
\end{tabular}
\end{center}

\paragraph{Confirmation.}
Pure Newton already achieves $\mathbb{E}[N_{\text{pre-basin}}] = O(1)$ on mvKS (max $7.8$ at $D = 1024$, scaling as $\sim D^{0.10}$, statistically indistinguishable from constant). ELS provides a $1.1$--$1.5\times$ improvement in pre-basin epoch count, never an asymptotic one. The univariate-style $\Omega(d/\log d)$ ELS speedup of Part~IX has \emph{no analogue} in the multivariate setting -- the multivariate Newton step is well-dimensioned by virtue of the matrix inversion.

\paragraph{Net conclusion.}
The mechanism that makes the univariate ELS rescue ABD (Lemma~3.1 of paper~101quater, $\tau^\star \asymp \sqrt d$) does \emph{not} drive the multivariate behaviour. Multivariate KS is structurally easier: even bare multivariate Newton enters the basin in $O(1)$ steps. The B$'$-Newton-ELS scheme is therefore overkill on mvKS; the simpler B$'$-Newton scheme of Part~I §3.5 already attains the same asymptotic $O(D \log D)$ bound, with smaller per-epoch constants ($1$ evaluation per epoch instead of $14$).

\section{Implications: $O(D \log D)$ Las Vegas bound on mvKS}\label{sec:final}

\begin{theorem}[Multivariate Las Vegas $O(D\log D + n^{3}\log D)$ on mvKS]\label{thm:mv-final}
Under the multivariate Kostlan--Smale ensemble of degree-$d$ systems on $\mathbb{C}^n$ (B\'ezout degree $D = d^n$), the multi-start B$'$-Newton-ELS algorithm of Definition~\ref{def:BNELS-mv} admits the empirically-supported bound
\begin{equation}\label{eq:mv-bound}
\mathbb{E}_{F \sim \mathrm{mvKS}}\!\bigl[\,\mathrm{cost}^{B'-N-\text{ELS}}_{\text{total}}(F)\,\bigr] \;=\; O\bigl(D \log D + n^{3}\log D\bigr),
\end{equation}
conditional on the multivariate analogues of paper~101bis Lemma~3.2 (Plancherel decorrelation on $\mathbb{C}^n$) and paper~101quater Lemma~3.1 (optimal-$\tau$ scaling, here $\tau^\star \asymp \sqrt D$).
\end{theorem}

\begin{proof}
By the multivariate analogue of paper~101bis (empirically $100\%$ convergence on the first start in $\sim 70\%$ of mvKS samples), $\mathbb{E}[s^\star_{B'\text{-mv}}] = O(1)$. By the regression \eqref{eq:fit-mv}, $\mathbb{E}[N_{\text{pre-basin}}] = O(1)$. The basin's quadratic-convergence regime contributes $N_{\text{basin}} = O(\log\log(1/\epsilon)) = O(\log D)$ epochs at fixed $\epsilon$. The per-epoch cost is $O(D + n^3)$ (one Jacobian solve + $14$ system evaluations). Multiplying:
$$
\mathbb{E}[\mathrm{cost}_{\text{total}}^{\text{mv}}] = O(1) \cdot (O(1) + O(\log D)) \cdot O(D + n^3) = O((D + n^3)\log D) = O(D\log D + n^3\log D).
$$
\end{proof}

\begin{remark}[Comparison with Beltr\'an--Pardo and Lairez]
For $n$ fixed and $d$ large, $D = d^n$ dominates $n^3$ and the bound becomes $O(D \log D)$. For comparison:
\begin{itemize}
\item Beltr\'an--Pardo~\cite{beltranpardo}: average polynomial cost in $D$, with explicit polynomial in $D$ (high-degree, via Smale $\mu$-theory + homotopy continuation).
\item Lairez~\cite{lairez}: deterministic average-case algorithm with cost polynomial in $D$.
\item Theorem~\ref{thm:mv-final}: $O(D \log D)$ for the specific B$'$-Newton-ELS scheme on mvKS, conditional on two structural lemmas.
\end{itemize}
The $O(D \log D)$ exponent in $D$ matches Lairez asymptotically (his exponent is $1$ + lower-order terms). We do \emph{not} claim a complexity-theoretic improvement; the contribution is to demonstrate that the simple Pandrosion-Newton-ELS algorithm reaches near-optimal scaling on mvKS without homotopy machinery.
\end{remark}

\begin{remark}[Honest scope]
The bound is empirical (regression on $9$ configurations, $108$ random systems). The two conditional lemmas (Plancherel decorrelation, optimal-$\tau$) need multivariate proofs, deferred to a future analytic companion. We do \emph{not} claim a complexity-theoretic resolution of Smale's 17th problem (already settled by~\cite{beltranpardo, lairez}).
\end{remark}

\section{Conclusion}\label{sec:conclusion}

The empirical $N_{\text{pre-basin}} = O(1)$ behaviour transfers from the univariate Part~IX setting to the multivariate Pandrosion framework of Part~I~\S3.5, with B$'$-Newton-ELS converging on mvKS up to B\'ezout degree $D = 1024$ in $\mathbb{E}[N_{\text{pre}}] \le 7$ epochs uniformly.

\paragraph{Honest revision: the mechanism differs from univariate.} The univariate ELS \emph{rescue} of Part~IX was driven by the systematic $\sqrt d$ undershooting of the Newton step (Lemma~3.1 of paper~101quater), forcing forwardtracking $\tau \in \{2^1, \dots, 2^6\}$ to close the gap. In the multivariate setting (this paper, \S\ref{sec:tau-diagnostic}), the observed optimal $\tau^\star_0$ is concentrated near $1$, \emph{not} $\sqrt D$: the multivariate Newton step is correctly dimensioned by the inverse Jacobian. Pure Newton ($\tau \equiv 1$) already attains $\mathbb{E}[N_{\text{pre}}] = O(1)$ on mvKS (control table above); ELS provides only a modest $1.1$--$1.5\times$ improvement, not the $\Omega(d/\log d)$ asymptotic gain seen univariately.

The Las~Vegas bound $\mathbb{E}[\mathrm{cost}_{\text{total}}^{\text{mv}}] = O(D \log D + n^3 \log D)$ on mvKS is the multivariate analogue of Part~IX's $O(d \log d)$ -- but the proof goes through whether one uses pure Newton or ELS. The $\log D$ gap to the trivial lower bound $\Omega(D)$ is the basin's $O(\log\log(1/\epsilon))$ for fixed precision; the $n^3 \log D$ term is the Jacobian-solve overhead.

\paragraph{Recommendation.} For multivariate root-finding on mvKS, use bare B$'$-Newton (the multivariate Pandrosion-with-dynamic-anchor of Part~I §3.5). The ELS extension is unnecessary in the multivariate regime and adds $14\times$ per-epoch overhead for at most $1.5\times$ epoch reduction.

We do \emph{not} claim a complexity-theoretic improvement upon Beltr\'an--Pardo~\cite{beltranpardo} or Lairez~\cite{lairez}.

\paragraph{Closure of the algorithmic series.} With Part~IX (univariate) and Part~IX-mv (multivariate, this paper), the algorithmic side of Universitas Pandrosion is now complete on both branches. Open analytic work: formal proofs of Plancherel decorrelation (univariate paper~101bis Lemma~3.2 and its multivariate analogue) and optimal-$\tau$ scaling (univariate paper~101quater Lemma~3.1 and its multivariate analogue).

\begin{thebibliography}{99}
\bibitem{besevic1} I. Besevic, \textit{A Pandrosion Operator for Derivative-Free Polynomial Root-Finding}. Universitas Pandrosion, Part~I (2026).
\bibitem{besevic7} I. Besevic, \textit{Universitas Pandrosion: Part~VII --- A derivative-free adaptive Pandrosion scheme with non-holomorphic fallback}. Preprint, 2026.
\bibitem{besevic8} I. Besevic, \textit{Universitas Pandrosion: Part~VIII --- The Geometric-Decreasing Start Strategy}. Preprint, 2026.
\bibitem{besevic9} I. Besevic, \textit{Universitas Pandrosion: Part~IX --- The Extended Line Search and the $O(d\log d)$ Algorithmic Bound}. Preprint, 2026.
\bibitem{besevic101bis} I. Besevic, \textit{Universitas Pandrosion: Paper~101bis --- The KS-Adaptive Start Radius and the Phase-Transition Elimination Theorem}. Preprint, 2026.
\bibitem{besevic101ter} I. Besevic, \textit{Universitas Pandrosion: Paper~101ter --- The Armijo Amortised Theorem}. Preprint, 2026.
\bibitem{besevic101quater} I. Besevic, \textit{Universitas Pandrosion: Paper~101quater --- The Extended Line Search Rescues ABD}. Preprint, 2026.
\bibitem{beltranpardo} C. Beltr\'an and L. M. Pardo, \textit{Smale's 17th problem: Average polynomial time to compute affine and projective solutions}. J. Amer. Math. Soc. \textbf{22} (2009), 363--385.
\bibitem{lairez} P. Lairez, \textit{A deterministic algorithm to compute approximate roots of polynomial systems in polynomial average time}. Found. Comput. Math. \textbf{17} (2017), 1265--1292.
\bibitem{schmidt1963} J. W. Schmidt, \textit{Eine \"Ubertragung der Regula Falsi auf Gleichungen in Banachr\"aumen}. Z. Angew. Math. Mech.\ \textbf{43} (1963), 97--110.
\bibitem{shubsmale} M. Shub and S. Smale, \textit{Complexity of B\'ezout's theorem I: Geometric aspects}. J. Amer. Math. Soc. \textbf{6} (1993), 459--501.
\bibitem{smale1986} S. Smale, \textit{Newton's method estimates from data at one point}. The Merging of Disciplines, Springer (1986), 185--196.
\bibitem{smale1998} S. Smale, \textit{Mathematical Problems for the Next Century}. The Mathematical Intelligencer, 20(2) (1998), 7--15.
\end{thebibliography}

\end{document}
